% !TEX root = ../Practica1_Hernandez_Andrea_Espino_Heriberto.tex
\usepackage[discipline=computer-science]{template/liverpoolthesis}
\usepackage[utf8]{inputenc}
\usepackage[T1]{fontenc}
\usepackage[spanish,es-nodecimaldot]{babel}

% Las referencias internas y las citas usan el mismo azul que los títulos.
\hypersetup{
  linkcolor = LiverpoolBlue,
  citecolor = LiverpoolBlue,
  urlcolor = LiverpoolBlue
}

% La marca de cada capítulo aparece en el encabezado, arriba de la regla.
\addto\captionsspanish{\renewcommand{\chaptername}{Ejercicio}}
\providecommand{\figureautorefname}{Figura}
\addto\captionsspanish{\renewcommand{\figureautorefname}{Figura}}

\graphicspath{{figuras/}}
\setlength{\headheight}{28pt}

% Inserta una captura de código y/o de ejecución.
% El tercer argumento es la etiqueta que se usa con \autoref{...}.
\newcommand{\incluirevidencia}[3]{%
  \begin{figure}[H]
    \centering
    \includegraphics[width=0.95\textwidth]{#1}
    \caption{#2}
    \label{#3}
  \end{figure}%
}

% ---------- Datos de la portada ----------
\title{Práctica de laboratorio 1: Creación de procesos e hilos en C y Java}
\author{%
  Heriberto Espino Montelongo\\
  ID: 175199\\[1em]
  Andrea Hernandez Huerta\\
  ID: 175523\\[1em]
  \text{Profesor:} Dr. Victor Giovanni Morales Murillo}
\faculty{Universidad de las Américas Puebla (UDLAP)}
\school{Licenciatura en Ciencia de Datos}
\degree{Curso: Programación Concurrente}
\date{Fecha de entrega: \today}
\subject{Creación y ejecución de procesos e hilos en C y Java}
